% 5. Projekt interfejsu
\section{Projekt interfejsu}\label{sec:projekt}

\begin{instrukcja}[zalecane]
Pokażcie drogę od makiety do gotowego ekranu i opiszcie decyzje wizualne z Wykładu 2: Ramy Graficzne
(stałe elementy na wszystkich ekranach), hierarchię wizualną (co dominuje na ekranie) i układ.
Makiety z K1 dajcie do załącznika A, a tutaj jedno porównanie makieta - ekran końcowy.
\end{instrukcja}

\subsection{Od makiety do ekranu}
\zrzut{makieta_lista.png}{Makieta listy z punktu kontrolnego K1}{0.8}
\zrzut{e2_lista.png}{Ekran listy w obecnej wersji (szeroki ekran, dwa panele)}{0.8}
Układ z makiety został zachowany: boczne menu, filtr w postaci przycisków, lista kart po lewej
i szczegóły wybranego terminu po prawej. Szarości zastąpił jeden kolor akcentu, a znaczniki typu
dostały ikony z zestawu Lucide obok tekstu. Przyciski „Edytuj termin” i „Usuń termin” z makiety
zostaną dodane razem z formularzem (L7).

\subsection{Ramy Graficzne i hierarchia wizualna}
Ekrany mają wspólny układ (\texttt{Layout.astro}) i wspólne style
(\texttt{global.css}): to samo menu w tym samym miejscu, ten sam krój, kolory, promienie i odstępy.
Na ekranie startowym dominuje karta najbliższego terminu: jest największa, ma grubszą ramkę, cień
i liczbę dni zapisaną cyfrą wielkości 4\,rem; dodatkowo wokół ramki krąży subtelny punkt światła.
Na liście dominuje zaznaczony termin i jego panel szczegółów, a na szczegółach blok z datą i liczbą dni.
Polecenie główne „Dodaj termin” to jedyny wypełniony przycisk w kolorze akcentu, z ikoną.

\subsection{Kolory i typografia}
\begin{table}[H]\centering
\caption{Kolory i kroje pisma}\label{tab:kolory}
\begin{tabularx}{\linewidth}{@{}p{3.4cm}p{2.6cm}L@{}}
\naglowektabeli \nt{Rola} & \nt{Wartość} & \nt{Gdzie używana} \\
kolor główny & \#2F3FD0 & przyciski główne, aktywna pozycja menu, linki, liczba dni \\
kolor tła & \#F6F7FB / \#FFFFFF & tło strony / karty i menu \\
kolor tekstu & \#1B1F2A / \#4A5163 & tekst podstawowy / pomocniczy \\
kolor błędu & \#B00020 & komunikaty o błędach w formularzu \\
kolor pilności & \#B3261E na \#FDECEA & znacznik „Za N dni” przy terminach pilnych \\
krój pisma & systemowy (\texttt{system-ui}, Segoe UI, Roboto) & cały interfejs; nie pobieramy fontów \\
\bottomrule
\end{tabularx}
\end{table}
